%% kapa: exam=1 round=20 year=2009 period=2 form=A subject=accounting subject_ko=회계학 seq=22 number=22
%% source: Q-Net 감정평가사 기출문제 (cst00309) artlSeq=5024178 file=extracted/1008/5024178/2009년 제20회 감정평가사 1차시험 기출문제/2교시.hwp
%% conversion: kordoc 4.24.51 -> tools/gen1.py (auto); review=proofread-sonnet; flags=none
\begin{question}{22}
(주)한국은 20×7년 1월 1일 우라늄광산을 ￦5,250,000에 취득하였다. 취득시 추정잔존가치는 ￦0이었고 추정매장량은 1,000,000톤이었다. 20×7년과 20×8년에 각각 80,000톤과 100,000톤의 채굴을 완료하였다. 20×8년도 채굴 완료 후 잔존매장량이 950,000톤인 것으로 추정되어, 20×8년도에 매장량의 추정을 변경하였다. 이러한 추정변경의 정당성은 입증되었다. 생산량비례법으로 상각하고 있다면, 20×8년도 상각액은 얼마인가?
\begin{choices}
  \item ￦420,000
  \item ￦442,105
  \item ￦460,000
  \item ￦508,421
  \item ￦525,000
\end{choices}
\end{question}
\answer{3}
